\documentclass[10pt,a4paper]{amsart}
\usepackage{amsmath,amssymb,amsthm,graphicx}
\usepackage{verbatim}
\newenvironment{reference}{\comment}{\endcomment}
%\newenvironment{reference}{}{}
\newenvironment{slogan}{\comment}{\endcomment}
\newenvironment{history}{\comment}{\endcomment}

% For commutative diagrams we use Xy-pic
\usepackage[all]{xy}

% We use 2cell for 2-commutative diagrams.
\xyoption{2cell}
% Theorem environments.
%
\theoremstyle{plain}
\newtheorem{theorem}[subsection]{Theorem}
\newtheorem{proposition}[subsection]{Proposition}
\newtheorem{lemma}[subsection]{Lemma}

\theoremstyle{definition}
\newtheorem{definition}[subsection]{Definition}
\newtheorem{example}[subsection]{Example}
\newtheorem{exercise}[subsection]{Exercise}
\newtheorem{situation}[subsection]{Situation}

\theoremstyle{remark}
\newtheorem{remark}[subsection]{Remark}
\newtheorem{remarks}[subsection]{Remarks}

\numberwithin{equation}{subsection}

% Macros
%
\def\lim{\mathop{\mathrm{lim}}\nolimits}
\def\colim{\mathop{\mathrm{colim}}\nolimits}
\def\Spec{\mathop{\mathrm{Spec}}}
\def\Hom{\mathop{\mathrm{Hom}}\nolimits}
\def\Ext{\mathop{\mathrm{Ext}}\nolimits}
\def\SheafHom{\mathop{\mathcal{H}\!\mathit{om}}\nolimits}
\def\SheafExt{\mathop{\mathcal{E}\!\mathit{xt}}\nolimits}
\def\Sch{\mathit{Sch}}
\def\Mor{\mathop{\mathrm{Mor}}\nolimits}
\def\Ob{\mathop{\mathrm{Ob}}\nolimits}
\def\Sh{\mathop{\mathit{Sh}}\nolimits}
\def\NL{\mathop{N\!L}\nolimits}
\def\CH{\mathop{\mathrm{CH}}\nolimits}
\def\proetale{{pro\text{-}\acute{e}tale}}
\def\etale{{\acute{e}tale}}
\def\QCoh{\mathit{QCoh}}
\def\Ker{\mathop{\mathrm{Ker}}}
\def\Im{\mathop{\mathrm{Im}}}
\def\Coker{\mathop{\mathrm{Coker}}}
\def\Coim{\mathop{\mathrm{Coim}}}

% Boxtimes
%
\DeclareMathSymbol{\boxtimes}{\mathbin}{AMSa}{"02}

%
% Macros for moduli stacks/spaces
%
\def\QCohstack{\mathcal{QC}\!\mathit{oh}}
\def\Cohstack{\mathcal{C}\!\mathit{oh}}
\def\Spacesstack{\mathcal{S}\!\mathit{paces}}
\def\Quotfunctor{\mathrm{Quot}}
\def\Hilbfunctor{\mathrm{Hilb}}
\def\Curvesstack{\mathcal{C}\!\mathit{urves}}
\def\Polarizedstack{\mathcal{P}\!\mathit{olarized}}
\def\Complexesstack{\mathcal{C}\!\mathit{omplexes}}
% \Pic is the operator that assigns to X its picard group, usage \Pic(X)
% \Picardstack_{X/B} denotes the Picard stack of X over B
% \Picardfunctor_{X/B} denotes the Picard functor of X over B
\def\Pic{\mathop{\mathrm{Pic}}\nolimits}
\def\Picardstack{\mathcal{P}\!\mathit{ic}}
\def\Picardfunctor{\mathrm{Pic}}
\def\Deformationcategory{\mathcal{D}\!\mathit{ef}}
\usepackage{cleveref}
\crefname{originallemma}{original lemma}{original lemmas}
\Crefname{originallemma}{Original lemma}{Original lemmas}
\makeatletter
\newlabel{algebra-lemma-completion-flat}{{00MB}{source}{}{sourcepointer-1}{}}
\newlabel{algebra-lemma-completion-flat@cref}{{[originallemma][0][]00MB}{[1][0][]source}{}{}{}}
\makeatother
\crefname{originallemma}{original lemma}{original lemmas}
\Crefname{originallemma}{Original lemma}{Original lemmas}
\makeatletter
\newlabel{algebra-lemma-etale-standard-smooth}{{00U9}{source}{}{sourcepointer-2}{}}
\newlabel{algebra-lemma-etale-standard-smooth@cref}{{[originallemma][0][]00U9}{[1][0][]source}{}{}{}}
\makeatother
\crefname{originalequation}{original equation}{original equations}
\Crefname{originalequation}{Original equation}{Original equations}
\makeatletter
\newlabel{equation-C-prime}{{0AL4}{source}{}{sourcepointer-3}{}}
\newlabel{equation-C-prime@cref}{{[originalequation][0][]0AL4}{[1][0][]source}{}{}{}}
\makeatother
\crefname{originallemma}{original lemma}{original lemmas}
\Crefname{originallemma}{Original lemma}{Original lemmas}
\makeatletter
\newlabel{formal-spaces-lemma-completion-in-sub}{{0ANB}{source}{}{sourcepointer-4}{}}
\newlabel{formal-spaces-lemma-completion-in-sub@cref}{{[originallemma][0][]0ANB}{[1][0][]source}{}{}{}}
\makeatother
\crefname{originallemma}{original lemma}{original lemmas}
\Crefname{originallemma}{Original lemma}{Original lemmas}
\makeatletter
\newlabel{formal-spaces-lemma-etale-surjective}{{0AN9}{source}{}{sourcepointer-5}{}}
\newlabel{formal-spaces-lemma-etale-surjective@cref}{{[originallemma][0][]0AN9}{[1][0][]source}{}{}{}}
\makeatother
\crefname{originalremark}{original remark}{original remarks}
\Crefname{originalremark}{Original remark}{Original remarks}
\makeatletter
\newlabel{formal-spaces-remark-variant-Noetherian}{{0ANI}{source}{}{sourcepointer-6}{}}
\newlabel{formal-spaces-remark-variant-Noetherian@cref}{{[originalremark][0][]0ANI}{[1][0][]source}{}{}{}}
\makeatother
\crefname{originalreference}{original reference}{original references}
\Crefname{originalreference}{Original reference}{Original references}
\makeatletter
\newlabel{lemma-NL-is-completion}{{0GAF}{source}{}{sourcepointer-7}{}}
\newlabel{lemma-NL-is-completion@cref}{{[originalreference][0][]0GAF}{[1][0][]source}{}{}{}}
\makeatother
\crefname{originalreference}{original reference}{original references}
\Crefname{originalreference}{Original reference}{Original references}
\makeatletter
\newlabel{lemma-exact-sequence-NL}{{00S2}{source}{}{sourcepointer-8}{}}
\newlabel{lemma-exact-sequence-NL@cref}{{[originalreference][0][]00S2}{[1][0][]source}{}{}{}}
\makeatother
\crefname{originalreference}{original reference}{original references}
\Crefname{originalreference}{Original reference}{Original references}
\makeatletter
\newlabel{lemma-finite-type}{{04U5}{source}{}{sourcepointer-9}{}}
\newlabel{lemma-finite-type@cref}{{[originalreference][0][]04U5}{[1][0][]source}{}{}{}}
\makeatother
\crefname{originalreference}{original reference}{original references}
\Crefname{originalreference}{Original reference}{Original references}
\makeatletter
\newlabel{lemma-transitive-lci-at-end}{{07D4}{source}{}{sourcepointer-10}{}}
\newlabel{lemma-transitive-lci-at-end@cref}{{[originalreference][0][]07D4}{[1][0][]source}{}{}{}}
\makeatother
\crefname{originallemma}{original lemma}{original lemmas}
\Crefname{originallemma}{Original lemma}{Original lemmas}
\makeatletter
\newlabel{more-algebra-lemma-base-change-RHom}{{0A6A}{source}{}{sourcepointer-11}{}}
\newlabel{more-algebra-lemma-base-change-RHom@cref}{{[originallemma][0][]0A6A}{[1][0][]source}{}{}{}}
\makeatother
\crefname{originallemma}{original lemma}{original lemmas}
\Crefname{originallemma}{Original lemma}{Original lemmas}
\makeatletter
\newlabel{more-algebra-lemma-base-change-property-ext-1-annihilated}{{0G9H}{source}{}{sourcepointer-12}{}}
\newlabel{more-algebra-lemma-base-change-property-ext-1-annihilated@cref}{{[originallemma][0][]0G9H}{[1][0][]source}{}{}{}}
\makeatother
\crefname{originallemma}{original lemma}{original lemmas}
\Crefname{originallemma}{Original lemma}{Original lemmas}
\makeatletter
\newlabel{more-algebra-lemma-two-term-base-change}{{0G9G}{source}{}{sourcepointer-13}{}}
\newlabel{more-algebra-lemma-two-term-base-change@cref}{{[originallemma][0][]0G9G}{[1][0][]source}{}{}{}}
\makeatother
\title{Rig-smooth is a local property on arrows of the category of adic Noetherian topological rings}
\author{The Stacks Project Authors}
\date{stacks-project:restricted.tex:local-snapshot-20261012; GNU FDL 1.2 or later}
\begin{document}
\maketitle
\noindent\textit{Editorial note.} All delivered mathematical text is the author's own, copied from the pinned Stacks Project chapter file restricted.tex; no statement or proof step was rewritten, reordered or removed.\par\medskip
\noindent\textit{Editorial note.} The counted proof is the author's complete argument: a preliminary base-change identification of completed naive cotangent complexes, followed by verification of the three locality conditions, with the descent step done by faithful flatness through the cited Ext base-change identity inside the delivered span\par\medskip
\noindent\textit{Editorial note.} The label lemma-rig-smooth-axioms names the counted statement, and the proof environment immediately following it in the source is delivered verbatim as its proof.\par\medskip
\noindent\textit{Editorial note.} The counted statement is carried by the author's own lemma environment in the chapter source.\par\medskip
\noindent\textit{Editorial note.} The proof of condition (2) applies the cited base-change identity with the author's parenthetical '(minor details omitted)'; the identity itself is a cited standard lemma (delivered as a navigation pointer), and the surrounding annihilation and faithful flatness arguments are complete in the delivered span.\par\medskip
\noindent\textit{Allowed standard prerequisites (editorial).}\begin{itemize}
\item The category WAdm\textasciicircum{}\{Noeth\} of adic Noetherian topological rings with continuous morphisms, and the local-property criterion of Formal Spaces for arrows
\item The completed naive cotangent complex with its functorial exact sequence, as developed in the same chapter
\item Base change and descent behaviour of the Ext annihilation characterizing rig-smoothness, from More on Algebra
\end{itemize}\medskip
\ifdefined\hypertarget\hypertarget{sourcepointer-1}{}\fi
\noindent\textit{Source reference.} Original lemma 00MB: sources/stacks-chow/tags.txt, source lines 817--817. Prerequisite pointer invoked at proof line(s) 4348; it names a standard fact of a sibling chapter that the counted argument uses, and every step of the counted proof built on it is retained verbatim in the delivered body.\par\medskip
\ifdefined\hypertarget\hypertarget{sourcepointer-2}{}\fi
\noindent\textit{Source reference.} Original lemma 00U9: sources/stacks-chow/tags.txt, source lines 1103--1103. Prerequisite pointer invoked at proof line(s) 4311; it names a standard fact of a sibling chapter that the counted argument uses, and every step of the counted proof built on it is retained verbatim in the delivered body.\par\medskip
\ifdefined\hypertarget\hypertarget{sourcepointer-3}{}\fi
\noindent\textit{Source reference.} Original equation 0AL4: sources/stacks-restricted/pinned-original-source.tex, source lines 326--326. Prerequisite pointer invoked at proof line(s) 4298; it names a standard fact of the same chapter that the counted argument uses, and every step of the counted proof built on it is retained verbatim in the delivered body.\par\medskip
\ifdefined\hypertarget\hypertarget{sourcepointer-4}{}\fi
\noindent\textit{Source reference.} Original lemma 0ANB: sources/stacks-chow/tags.txt, source lines 13434--13434. Prerequisite pointer invoked at proof line(s) 4278; it names a standard fact of a sibling chapter that the counted argument uses, and every step of the counted proof built on it is retained verbatim in the delivered body.\par\medskip
\ifdefined\hypertarget\hypertarget{sourcepointer-5}{}\fi
\noindent\textit{Source reference.} Original lemma 0AN9: sources/stacks-chow/tags.txt, source lines 13432--13432. Prerequisite pointer invoked at proof line(s) 4359; it names a standard fact of a sibling chapter that the counted argument uses, and every step of the counted proof built on it is retained verbatim in the delivered body.\par\medskip
\ifdefined\hypertarget\hypertarget{sourcepointer-6}{}\fi
\noindent\textit{Source reference.} Original remark 0ANI: sources/stacks-chow/tags.txt, source lines 13441--13441. Prerequisite pointer invoked at proof line(s) 4256; it names a standard fact of a sibling chapter that the counted argument uses, and every step of the counted proof built on it is retained verbatim in the delivered body.\par\medskip
\ifdefined\hypertarget\hypertarget{sourcepointer-7}{}\fi
\noindent\textit{Source reference.} Original reference 0GAF: sources/stacks-restricted/pinned-original-source.tex, source lines 560--560. Prerequisite pointer invoked at proof line(s) 4302; it names a standard fact of the same chapter that the counted argument uses, and every step of the counted proof built on it is retained verbatim in the delivered body.\par\medskip
\ifdefined\hypertarget\hypertarget{sourcepointer-8}{}\fi
\noindent\textit{Source reference.} Original reference 00S2: sources/stacks-restricted/pinned-original-source.tex, source lines 725--725. Prerequisite pointer invoked at proof line(s) 4303; it names a standard fact of the same chapter that the counted argument uses, and every step of the counted proof built on it is retained verbatim in the delivered body.\par\medskip
\ifdefined\hypertarget\hypertarget{sourcepointer-9}{}\fi
\noindent\textit{Source reference.} Original reference 04U5: sources/stacks-restricted/pinned-original-source.tex, source lines 2439--2439. Prerequisite pointer invoked at proof line(s) 4287; it names a standard fact of the same chapter that the counted argument uses, and every step of the counted proof built on it is retained verbatim in the delivered body.\par\medskip
\ifdefined\hypertarget\hypertarget{sourcepointer-10}{}\fi
\noindent\textit{Source reference.} Original reference 07D4: sources/stacks-restricted/pinned-original-source.tex, source lines 780--780. Prerequisite pointer invoked at proof line(s) 4314; it names a standard fact of the same chapter that the counted argument uses, and every step of the counted proof built on it is retained verbatim in the delivered body.\par\medskip
\ifdefined\hypertarget\hypertarget{sourcepointer-11}{}\fi
\noindent\textit{Source reference.} Original lemma 0A6A: sources/stacks-chow/tags.txt, source lines 12838--12838. Prerequisite pointer invoked at proof line(s) 4355; it names a standard fact of a sibling chapter that the counted argument uses, and every step of the counted proof built on it is retained verbatim in the delivered body.\par\medskip
\ifdefined\hypertarget\hypertarget{sourcepointer-12}{}\fi
\noindent\textit{Source reference.} Original lemma 0G9H: sources/stacks-chow/tags.txt, source lines 20302--20302. Prerequisite pointer invoked at proof line(s) 4335; it names a standard fact of a sibling chapter that the counted argument uses, and every step of the counted proof built on it is retained verbatim in the delivered body.\par\medskip
\ifdefined\hypertarget\hypertarget{sourcepointer-13}{}\fi
\noindent\textit{Source reference.} Original lemma 0G9G: sources/stacks-chow/tags.txt, source lines 20301--20301. Prerequisite pointer invoked at proof line(s) 4334; it names a standard fact of a sibling chapter that the counted argument uses, and every step of the counted proof built on it is retained verbatim in the delivered body.\par\medskip
\begin{lemma}
\label{lemma-rig-smooth-axioms}
The property $P(\varphi)=$``$\varphi$ is rig-smooth'' on arrows
of $\textit{WAdm}^{Noeth}$ is a local property as defined in
Formal Spaces, Remark \ref{formal-spaces-remark-variant-Noetherian}.
\end{lemma}

\begin{proof}
Let us recall what the statement signifies. First, 
$\textit{WAdm}^{Noeth}$ is the category whose objects are
adic Noetherian topological rings and whose morphisms are
continuous ring homomorphisms. Consider a commutative diagram
$$
\xymatrix{
B \ar[r] & (B')^\wedge \\
A \ar[r] \ar[u]^\varphi & (A')^\wedge \ar[u]_{\varphi'}
}
$$
satisfying the following conditions:
$A$ and $B$ are adic Noetherian topological rings,
$A \to A'$ and $B \to B'$ are \'etale ring maps,
$(A')^\wedge = \lim A'/I^nA'$ for some ideal of definition $I \subset A$,
$(B')^\wedge = \lim B'/J^nB'$ for some ideal of definition $J \subset B$, and
$\varphi : A \to B$ and $\varphi' : (A')^\wedge \to (B')^\wedge$
are continuous. Note that $(A')^\wedge$ and $(B')^\wedge$ are
adic Noetherian topological rings by
Formal Spaces, Lemma \ref{formal-spaces-lemma-completion-in-sub}.
We have to show
\begin{enumerate}
\item $\varphi$ is rig-smooth $\Rightarrow \varphi'$ is rig-smooth,
\item if $B \to B'$ faithfully flat, then $\varphi'$ is rig-smooth
$\Rightarrow \varphi$ is rig-smooth, and
\item if $A \to B_i$ is rig-smooth for $i = 1, \ldots, n$, then
$A \to \prod_{i = 1, \ldots, n} B_i$ is rig-smooth.
\end{enumerate}
The equivalent conditions of Lemma \ref{lemma-finite-type} satisfy
conditions (1), (2), and (3).
Thus in verifying (1), (2), and (3) for the property
``rig-smooth'' we may already assume our ring maps satisfy
the equivalent conditions of Lemma \ref{lemma-finite-type}
in each case.

\medskip\noindent
Pick an ideal of definition $I \subset A$. By the remarks above
the topology on each ring in the diagram is the $I$-adic topology
and $B$, $(A')^\wedge$, and $(B')^\wedge$ are in the category
(\ref{equation-C-prime}) for $(A, I)$.
Since $A \to A'$ and $B \to B'$ are \'etale the complexes
$\NL_{A'/A}$ and $\NL_{B'/B}$ are zero and hence
$\NL_{(A')^\wedge/A}^\wedge$ and $\NL_{(B')^\wedge/B}^\wedge$
are zero by Lemma \ref{lemma-NL-is-completion}.
Applying Lemma \ref{lemma-exact-sequence-NL} to
$A \to (A')^\wedge \to (B')^\wedge$ we get isomorphisms
$$
H^i(\NL_{(B')^\wedge/(A')^\wedge}^\wedge) \to H^i(\NL_{(B')^\wedge/A}^\wedge)
$$
Thus $\NL_{(B')^\wedge/A}^\wedge \to \NL_{(B')^\wedge/(A')^\wedge}$
is a quasi-isomorphism. The ring maps $B/I^nB \to B'/I^nB'$ are \'etale
and hence are local complete intersections
(Algebra, Lemma \ref{algebra-lemma-etale-standard-smooth}).
Hence we may apply
Lemmas \ref{lemma-exact-sequence-NL} and
\ref{lemma-transitive-lci-at-end} to
$A \to B \to (B')^\wedge$ and we get isomorphisms
$$
H^i(\NL_{B/A}^\wedge \otimes_B (B')^\wedge) \to
H^i(\NL_{(B')^\wedge/A}^\wedge)
$$
We conclude that
$\NL_{B/A}^\wedge \otimes_B (B')^\wedge \to \NL_{(B')^\wedge/A}^\wedge$
is a quasi-isomorphism. Combining these two observations we obtain that
$$
\NL_{(B')^\wedge/(A')^\wedge}^\wedge \cong
\NL_{B/A}^\wedge \otimes_B (B')^\wedge
$$
in $D((B')^\wedge)$.
With these preparations out of the way we can start the actual proof.

\medskip\noindent
Proof of (1). Assume $\varphi$ is rig-smooth. Then there exists a $c \geq 0$
such that $\Ext^1_B(\NL_{B/A}^\wedge, N)$ is annihilated by $I^c$
for every $B$-module $N$. By
More on Algebra, Lemmas \ref{more-algebra-lemma-two-term-base-change} and
\ref{more-algebra-lemma-base-change-property-ext-1-annihilated}
this property is preserved under base change by $B \to (B')^\wedge$.
Hence $\Ext^1_{(B')^\wedge}(\NL_{(B')^\wedge/(A')^\wedge}^\wedge, N)$
is annihilated by $I^c(A')^\wedge$ for all $(B')^\wedge$-modules $N$
which tells us that $\varphi'$ is rig-smooth.
This proves (1).

\medskip\noindent
To prove (2) assume $B \to B'$ is faithfully flat and that $\varphi'$
is rig-smooth. Then there exists a $c \geq 0$ such that
$\Ext^1_{(B')^\wedge}(\NL_{(B')^\wedge/(A')^\wedge}^\wedge, N')$
is annihilated by $I^c(B')^\wedge$ for every $(B')^\wedge$-module $N'$.
The composition $B \to B' \to (B')^\wedge$ is flat
(Algebra, Lemma \ref{algebra-lemma-completion-flat})
hence for any $B$-module $N$ we have
$$
\Ext^1_B(\NL_{B/A}^\wedge, N) \otimes_B (B')^\wedge =
\Ext^1_{(B')^\wedge}(\NL_{B/A}^\wedge \otimes_B (B')^\wedge,
N \otimes_B (B')^\wedge)
$$
by More on Algebra, Lemma \ref{more-algebra-lemma-base-change-RHom} part (3)
(minor details omitted). Thus we see that this module is annihilated
by $I^c$. However, $B \to (B')^\wedge$ is actually faithfully flat
by our assumption that $B \to B'$ is faithfully flat (Formal Spaces, Lemma
\ref{formal-spaces-lemma-etale-surjective}). Thus we conclude that
$\Ext^1_B(\NL_{B/A}^\wedge, N)$ is annihilated by $I^c$.
Hence $\varphi$ is rig-smooth. This proves (2).

\medskip\noindent
To prove (3), setting $B = \prod_{i = 1, \ldots, n} B_i$
we just observe that $\NL_{B/A}^\wedge$ is the direct
sum of the complexes $\NL_{B_i/A}^\wedge$ viewed as complexes
of $B$-modules.
\end{proof}

\end{document}
